\documentclass[11pt,a4paper]{article}
\usepackage{fontspec,amsmath,amssymb,amsthm,microtype}
\setmainfont{lmroman10-regular.otf}[BoldFont=lmroman10-bold.otf,ItalicFont=lmroman10-italic.otf,BoldItalicFont=lmroman10-bolditalic.otf]
\usepackage{xcolor}
\usepackage[most]{tcolorbox}
\usepackage{mathrsfs}
\providecommand{\lt}{<}
\providecommand{\gt}{>}
\setmonofont{DejaVuSansMono.ttf}[BoldFont=DejaVuSansMono-Bold.ttf,ItalicFont=DejaVuSansMono-Oblique.ttf,BoldItalicFont=DejaVuSansMono-BoldOblique.ttf,Scale=0.86]
\newfontfamily\leanfallback{STIXTwoMath-Regular.otf}[Scale=0.86]
\definecolor{leanrule}{gray}{0.62}
\definecolor{leanback}{gray}{0.962}
\definecolor{leantitle}{gray}{0.22}
\definecolor{leanmuted}{gray}{0.38}
\newtcolorbox{leanbox}[1]{enhanced,breakable,colback=leanback,colframe=leanrule,
  boxrule=0.4pt,arc=0.6mm,left=2.4mm,right=2.4mm,top=1.8mm,bottom=1.8mm,
  before skip=0.8em,after skip=1em,fonttitle=\small\bfseries,coltitle=leantitle,
  colbacktitle=leanback,titlerule=0.3pt,title={#1}}
\newlength{\leancharwidth}
\newenvironment{leancode}{\par\smallskip\begingroup\ttfamily\footnotesize\raggedright
  \settowidth{\leancharwidth}{M}\setlength{\parskip}{0pt}\setlength{\parindent}{0pt}
  \frenchspacing}{\par\endgroup\smallskip}
\newcommand{\leanline}[2]{\hangindent=\dimexpr#1\leancharwidth+4\leancharwidth\relax
  \hangafter=1\noindent\hspace*{#1\leancharwidth}#2\par}
\newcommand{\leansource}[1]{\par{\footnotesize\color{leanmuted}#1}\par}
\newcommand{\leanlabel}[1]{\par\smallskip\noindent{\small\bfseries #1}\par}
\usepackage[margin=25mm]{geometry}
\usepackage[unicode,hidelinks]{hyperref}
\hypersetup{pdfauthor={LangevinLean Project}}
\setlength{\emergencystretch}{3em}
\setlength{\parskip}{0.4em}
\setlength{\parindent}{0pt}
\allowdisplaybreaks[2]
\raggedbottom
\title{\texorpdfstring{State-measurable random variables are actual evaluated observables}{State-measurable random variables are actual evaluated observables}}
\author{LangevinLean Project}
\date{}
\begin{document}\maketitle
Draft mathematical exposition; these notes may contain mistakes.


Doob–Dynkin factorization identifies every square-integrable function of a state with an observable in the \(L^2\) space of its actual law. On this range, evaluation followed by its adjoint is the identity, identifying state-measurable conditional means with actual adjoint compression.


\section{\texorpdfstring{Actual square-integrable Doob–Dynkin factorization}{Actual square-integrable Doob–Dynkin factorization}}


\textbf{Claim.} If \(X_*Q=\mu\) and \(Z\in L^2(Q)\) is measurable from \(X\) alone, then \(Z=f(X)\) for some actual \(f\in L^2(\mu)\).



\textbf{Proof.} Choose a state-measurable representative of \(Z\). Measurable factorization writes it as \(g\circ X\). Since this is almost surely the square-integrable original variable, its square norm is integrable. The exact pushforward identity transfers that integrability to \(g\) under \(\mu\). Its \(L^2(\mu)\) equivalence class evaluates back to \(Z\).

\begin{leanbox}{Lean: \texttt{observable\_\allowbreak{}range\_\allowbreak{}of\_\allowbreak{}state\_\allowbreak{}measurable}}
\leansource{\texttt{observable\_\allowbreak{}range\_\allowbreak{}of\_\allowbreak{}state\_\allowbreak{}measurable} \quad WassersteinMixing/\allowbreak{}ObservableRange.lean, lines 27–43}
\begin{leancode}
\leanline{0}{theorem\ observable\_range\_of\_state\_measurable\ (Z\ :\ Lp\ F\ 2\ Q)}
\leanline{4}{(hZ\ :\ AEStronglyMeasurable[mA.comap\ X]\ Z\ Q)\ :}
\leanline{4}{∃\ f\ :\ Lp\ F\ 2\ μ,\ Lp.compMeasurePreservingₗᵢ\ ℝ\ X\ hX\ f\ =\ Z}
\end{leancode}
\leanlabel{How the Lean proof follows the argument}
\texttt{observable\_\allowbreak{}range\_\allowbreak{}of\_\allowbreak{}state\_\allowbreak{}measurable} uses measurable composition factorization on the chosen strongly measurable representative. It transfers \texttt{Mem\allowbreak{}Lp} through the actual measure-preserving state map before constructing the observable class and comparing pullback representatives.
\end{leanbox}

\section{\texorpdfstring{Identifying the true conditional mean with compression}{Identifying the true conditional mean with compression}}


\textbf{Claim.} For an evaluation isometry \(J\), \(JJ^*Z=Z\) on its range. If all evaluations are past-measurable and \(\mathbb E[Y\mid\mathcal F]\) belongs to that range, then

\[\mathbb E[Y\mid\mathcal F]=J(J^*Y).\]

\textbf{Proof.} The isometry identity \(J^*J=I\) proves the first claim. Pair the conditional expectation against arbitrary evaluated observables. Adaptation lets the conditional projection be removed from that pairing, yielding \(J^*\mathbb E[Y\mid\mathcal F]=J^*Y\). Apply \(J\) and use the first claim for the known range member.

\begin{leanbox}{Lean: \texttt{evaluation\_\allowbreak{}projection\_\allowbreak{}eq\_\allowbreak{}of\_\allowbreak{}mem\_\allowbreak{}range}, \texttt{conditional\_\allowbreak{}eq\_\allowbreak{}evaluation\_\allowbreak{}adjoint}}
\leansource{\texttt{evaluation\_\allowbreak{}projection\_\allowbreak{}eq\_\allowbreak{}of\_\allowbreak{}mem\_\allowbreak{}range} \quad WassersteinMixing/\allowbreak{}ObservableRange.lean, lines 54–59}
\begin{leancode}
\leanline{0}{theorem\ evaluation\_projection\_eq\_of\_mem\_range\ (Z\ :\ Lp\ F\ 2\ Q)\ (hZ\ :\ Z\ ∈\ Set.range\ J)\ :}
\leanline{4}{J\ (J.toContinuousLinearMap.adjoint\ Z)\ =\ Z}
\end{leancode}
\leansource{\texttt{conditional\_\allowbreak{}eq\_\allowbreak{}evaluation\_\allowbreak{}adjoint} \quad WassersteinMixing/\allowbreak{}ObservableRange.lean, lines 63–77}
\begin{leancode}
\leanline{0}{theorem\ conditional\_eq\_evaluation\_adjoint\ \{m\ :\ MeasurableSpace\ Ω\}\ (hm\ :\ m\ ≤\ mΩ)}
\leanline{4}{(Y\ :\ Lp\ F\ 2\ Q)}
\leanline{4}{(hJ\ :\ ∀\ f,\ AEStronglyMeasurable[m]\ (J\ f)\ Q)}
\leanline{4}{(hZ\ :\ (condExpL2\ F\ ℝ\ hm\ Y\ :\ Lp\ F\ 2\ Q)\ ∈\ Set.range\ J)\ :}
\leanline{4}{(condExpL2\ F\ ℝ\ hm\ Y\ :\ Lp\ F\ 2\ Q)\ =\ J\ (J.toContinuousLinearMap.adjoint\ Y)}
\end{leancode}
\leanlabel{How the Lean proof follows the argument}
\texttt{evaluation\_\allowbreak{}projection\_\allowbreak{}eq\_\allowbreak{}of\_\allowbreak{}mem\_\allowbreak{}range} cancels the isometry adjoint on an explicit range witness. \texttt{conditional\_\allowbreak{}eq\_\allowbreak{}evaluation\_\allowbreak{}adjoint} proves equality of extracted adjoint vectors by testing against every observable and using the conditional-expectation pairing identity.
\end{leanbox}

\end{document}
